\section{Discussion}

This article asked how much of the working day in the three largest community service occupations is spent on tasks that current AI can realistically speed up. The analysis shows that directly exposed tasks make up a small part of that day. Social and Human Service Assistants and Child, Family, and School Social Workers each spend less than an hour a day on them. A larger share of the day falls on partially exposed tasks, which AI could speed up only with software built for the purpose. This pattern holds all three occupations, and under both human and GPT-4 ratings, where a larger portion of the workday is partially exposed rather than directly exposed.

The measured exposure of community service work to AI depends on how exposure is counted. When tasks are counted equally, 32 percent of the work of Social and Human Service Assistants is directly exposed to AI. When tasks are weighted by time, the figure is just 6 percent of their day. The tasks in this occupation that an LLM could accelerate without additional software, such as referring clients to other agencies or submitting reports to a supervisor, occupy little of the working day. This result is consistent with Hatgis-Kessell et al. (2026), who find across the economy that time weighting reduces the contribution of short, transactional tasks. Task-based measures therefore yield a more favorable estimate of what AI can immediately offer these occupations than time-based measures do.

Most of the exposed time in these occupations is partially exposed, and it is concentrated in a small number of long tasks. For Social and Human Service Assistants, a single task, developing and implementing care plans, accounts for almost two-thirds of partially exposed time. Under the definition used by Eloundou et al. (2024), such tasks could be accelerated only if additional software were built on an LLM. The larger share of potential time savings in these occupations is therefore not accessible through general-purpose AI tools. It is conditional on software that would need to be developed, funded, and maintained.

Observed use was recorded to only four tasks across all three occupations. None of the four tasks were rated as directly exposed. This should not be interpreted as evidence that this workforce does not use AI. The observed use metric from Massenkoff and McCrory (2026) covers only one platform (Claude), and 64 percent of surveyed social workers report that they are using AI in their work (Borah et al. 2026). It does indicate that observed use has not occurred on the tasks the tasks that were rated as the most theoretically exposed.

\setcounter{subsection}{1} % keeps the draft numbering (5.2)

\subsection{Implications for policy and practice.}

\emph{For nonprofit managers.} These results suggest that general-purpose AI tools, used without additional software, are likely to return a limited amount of frontline staff time. Under the human ratings, the minimum saving on directly exposed tasks is about 12 minutes a day for a Social and Human Service Assistant and about 18 minutes for a Child, Family, and School Social Worker. Savings of this size may be worthwhile, but they are unlikely to offset the staffing shortfalls described in Section 2. Managers who expect AI to relieve frontline workloads should distinguish between tasks that current tools can accelerate, which are brief, and the longer tasks, such as care planning and case recording, that would require software designed to augment the LLM.

\emph{For funders.} Most of the working time that AI could affect in these occupations is partially exposed, and it is concentrated in a small number of tasks that are common across organizations. Few local community service organizations have the capacity to commission software or get specialized training for AI adoption (Smith Arrillaga, Grundhoefer, and Im 2025). Custom-built software with AI features are not common in this sector: just 5 percent of social workers describe using such software for case management (Borah et al. 2026). Building dedicated software for time-consuming tasks such as care planning may therefore return more staff time than support for individual licenses to general-purpose AI chat products.

\emph{For policymakers}. The single largest partially exposed tasks in the frontline occupations (care planning and case recording) likely can only be accelerated only with dedicated software built to extend the LLM. Agencies that require grantees to use specific data systems should look for ways to support AI features inside of those systems. Guidance on the use of client information in AI tools would address the barriers social workers cite most often, ethics and data privacy (Borah et al. 2026).

\subsection{Limitations}

The findings depend on three published measures, each with limitations. The exposure ratings date from early 2023 and are subjective: human and GPT-4 ratings agree on about half of the tasks examined, and the result for managers differs widely between them (Section 4.5). The time estimates are generated by a language model, not observed, and assign a large share of the day to single tasks. The measure of observed use covers one AI product and records a task only above a minimum volume, so it understates use to an unknown degree.

The scope of the analysis is also limited. Employment data are available only for the industry group as a whole, which includes for-profit employers and organizations outside the population of interest. The three occupations examined account for 27.6 percent of the industry's jobs, and the results should not be generalized to the sector. Estimates of total worker hours treat every job as full-time and exclude volunteers, who are common in food and shelter organizations.
